\documentclass[10pt]{amsart}
\usepackage[a4paper,margin=23mm]{geometry}
\usepackage{amsmath,amssymb,amsthm,xfrac,hyperref,xspace}
\usepackage[english]{babel}
\usepackage[scr=rsfs,cal=euler]{mathalfa}
\newcommand{\resp}{resp.}
\newcommand{\cf}{cf.}
\newcommand{\ie}{i.e.}
\newcommand{\eg}{e.g.}
\newcommand{\loccit}{loc.\ cit.}
\newcommand{\skpt}{\par\smallskip\noindent}
\let\Subsection\subsection
\let\Subsubsection\subsubsection
\emergencystretch=3em
\multlinegap0pt

\hyphenation{assump-tion believe between define defined defines dynam-ics ener-gy essen-tial-ly inte-gra-ble momen-tum param-e-ter}

\let\dpl\displaystyle
\let\ts\textstyle

\newcommand\mto{\mathchoice{\longmapsto}{\mapsto}{\mapsto}{\mapsto}}

\newcommand{\Psfrac}[2]{(\sfrac{#1}{#2})}
\newcommand{\psfrac}[2]{\sfrac{(#1)}{#2}}
\newcommand{\spfrac}[2]{\sfrac{#1}{(#2)}}
\newcommand{\pspfrac}[2]{\sfrac{(#1)}{(#2)}}

\DeclareMathOperator{\rw}{w}

\theoremstyle{plain}
\newtheorem{theorem}{Theorem}[section]
\newtheorem{proposition}[theorem]{Proposition}
\newtheorem{corollary}[theorem]{Corollary}
\newtheorem{lemma}[theorem]{Lemma}

\theoremstyle{definition}
\newtheorem{definition}[theorem]{Definition}
\newtheorem{notation}[theorem]{Notation}
\newtheorem{remark}[theorem]{Remark}

\DeclareMathOperator{\e}{e}
\DeclareMathOperator{\hi}{hi}
\DeclareMathOperator{\Id}{Id}
\DeclareMathOperator{\imm}{Im}\let\im\imm
\DeclareMathOperator{\lo}{lo}
\DeclareMathOperator{\meas}{meas}
\DeclareMathOperator{\ree}{Re}\let\re\ree
\DeclareMathOperator{\supp}{supp}
\DeclareMathOperator{\vol}{Vol}

\newcommand{\nb}{{\boldsymbol n}}
\newcommand{\kb}{{\boldsymbol k}}
\newcommand{\lb}{{\boldsymbol \ell}}
\newcommand{\mb}{{\boldsymbol m}}
\newcommand{\bb}{{\boldsymbol b}}
\newcommand{\hb}{{\boldsymbol h}}

\newcommand{\Nb}{{\boldsymbol{\mathfrak{N}}}}

\newcommand{\N}{\mathbb{N}}
\newcommand{\Z}{\mathbb{Z}}
\newcommand{\Q}{\mathbb{Q}}
\newcommand{\R}{\mathbb{R}}
\newcommand{\C}{\mathbb{C}}
\newcommand{\T}{\mathbb{T}}

\newcommand{\1}{\mathbf{1}}

\newcommand{\Clo}{\mathscr{C}_{\lo}}
\newcommand{\Chigh}{\mathscr{C}_{\hi}}

\newcommand{\Ysup}[1]{Y_{#1}^{\mathrm{sup}}}
\newcommand{\Ylip}[1]{Y_{#1}^{\mathrm{lip}}}

\let\epsilon\varepsilon

\begin{document}
\title{A priori estimates for the high super-actions}
\author{Joackim Bernier \and Nicolas Camps}
\maketitle
\section{Source equation and notation}
\subsection{Context}
We consider the nonlinear Schrödinger equation
\begin{equation}
\label{eq:nls}
\tag{NLS}
i\partial_{t}u + \mathrm{div} \, G \nabla u = f(|u|^{2})u,\quad (t,x)\in\R\times\mathbb{T}^{d},
\end{equation}
where $G$ is a real symmetric $d\times d$ positive matrix and $f:\R\to\R$ is a $C^\infty$ function satisfying $f'(0)\neq 0$ (to ensure the existence of a cubic nonlinear term in the equation). Without loss of generality, we~assume that $f(0)=0$.

This is a convenient way of rewriting nonlinear Schrödinger equations on rescaled tori. Indeed, $d$-dimensional flat tori writes $\mathbb{T}_{\mathscr{L}}^d= \R^d/\mathscr{L}$ where $\mathscr{L}=: V \mathbb{Z}^d$ is a lattice of $\R^d$ ($V$ being a $d\times d$ real invertible matrix) and so by setting $\operatorname{G} = V^{-1} \, ^{t}(\operatorname{V}^{-1})$ and by applying a linear change of coordinate $x\mto Vx$, \eqref{eq:nls} is equivalent to
\begin{equation*}
i\partial_tu + \Delta u = f(|u|^2)u,\quad (t,x)\in \R\times\T_{\mathscr{L}}^d.
\end{equation*}

By considering small solutions, \eqref{eq:nls} can be seen as a perturbation of the linear integrable system
\begin{equation}
\label{eq:Schrolin}
i\partial_{t}u + \mathrm{div} \, G \nabla u = 0
\end{equation}
whose actions $I_n(u) = |u_n|^2$ are the square modulus of the Fourier coefficients
\[
u_{n} = \frac{1}{(2\pi)^d}\int_{\T^{d}}\e^{-in\cdot x}u(x)\, \mathrm{d}x.
\]
\[
\lambda_n^2:= g(n,n), \quad \mathrm{where} \quad g(a,b):={} ^t\!a\,G\, b.
\]

To continue the discussion of the proof, we~need to introduce the Hamiltonian structure of \eqref{eq:nls}. More precisely, \eqref{eq:nls} rewrites
\begin{equation}
\label{eq:hamNLS}
i\partial_t u = \nabla H(u) \quad \mathrm{with} \quad H(u) = Z_2(u)+ \frac12 \int_{\mathbb{T}^d} F(|u(x)|^2) \, \mathrm{d}x,
\end{equation}
where $F$ is the primitive of $f$ vanishing at the origin and
\[
Z_2(u) = \frac12 \sum_{n\in \mathbb{Z}^d} \lambda_n^2 |u_n|^2.
\]
The classical notations (like $\nabla$) are defined in Section \ref{sec:notations} below.
\setcounter{subsection}{3}
\Subsection{Notations and functional setting}
\label{sec:notations}
\subsubsection{Functional setting}
\label{sec:fun}
We equip the torus $\mathbb{T}^d = \mathbb{R}^d / 2\pi \mathbb{Z}^d$ with the normalized Lebesgue measure $(2\pi)^{-d}\mathrm{d}x$. Therefore, the Lebesgue norms $\|\cdot\|_{L^p}$, $1\leq p<\infty$, are defined by density through the formula\vspace*{-3pt}
\[
\forall u \in C^0(\mathbb{T}^d), \quad \| u\|_{L^p}^p:= (2\pi)^{-d}\int_{\mathbb{T}^d} |u(x)|^p \mathrm{d}x.
\]
We identity each distribution $u$ on $\mathbb{T}^d$ with the sequence of its Fourier coefficients defined by density by\vspace*{-3pt}
\[
u_k = (2\pi)^{-d} \int_{\mathbb{T}^d} u(x) e^{-ikx} \mathrm{d}x.
\]
Given $1\leq p < \infty$ and $s\in \mathbb{R}$ we set
\[
\ts
\ell^{p}_s(\mathbb{Z}^d):= \bigl\{ u\in \mathbb{C}^{\mathbb{Z}^d} \mid \|u\|_{\ell^p_s}^p := \sum_{k\in \mathbb{Z}^d} \langle k \rangle^{ps} |u_k|^p <\infty \bigr\}.
\]
As usual, we~set
\[
H^s(\mathbb{T}^d) = h^s(\mathbb{Z}^d) = \ell^2_s(\mathbb{Z}^d) \quad \mathrm{and} \quad \ell^p(\mathbb{Z}^d) := \ell^p_0(\mathbb{Z}^d).
\]
Note that with these conventions, the Fourier--Plancherel isometry writes
\[
\forall u \in L^2(\mathbb{T}^d), \quad \| u\|_{L^2}^2 = \|u\|_{\ell^2}^2.
\]
Throughout this paper, we~consider $L^2$ and $\ell^2$ as real Hilbert spaces equipped with the scalar products
\[
\forall u,v\in L^2(\mathbb{T}^d), \quad (u,v)_{L^2} :=(2\pi)^{-d} \re \int_{\mathbb{T}^d} u(x) \overline{v(x)} \mathrm{d}x = \re \sum_{k\in \mathbb{Z}^d} u_k \overline{v_k} =: (u,v)_{\ell^2}.
\]
Note that if $s \in \mathbb{N}$ is a nonnegative integer, we~have
\[
\forall u \in H^s(\mathbb{T}), \quad \|u\|_{H^s}^2 = (2\pi)^{-d} \int_{\mathbb{T}^d} |\partial_x^s u(x)|^2 \mathrm{d}x.
\]

\subsubsection{Differential calculus and Poisson brackets}

Given $p\in[1,\infty)$, $s\in \mathbb{R}$, $\mathcal{U}$ an open subset of $\ell^p_s(\mathbb{Z}^d) $, a smooth function $P :\mathcal{U} \to \mathbb{R}$ and $u\in \ell^p_s(\mathbb{Z}^d)$, its gradient $\nabla P(u)$ is the unique element of $\ell^{p'}_{-s}(\mathbb{Z}^d)$ satisfying
\[
\forall v\in \ell^p_s(\mathbb{Z}^d), \ (\nabla P(u),v)_{L^2} = \mathrm{d}P(u)(v).
\]
It can be checked that
\begin{equation}
\label{eq:formula_grad}
\forall k \in \mathbb{Z}^d, \quad (\nabla P(u))_k = 2\partial_{\overline{u_k}} P(u).
\end{equation}
We equip $L^2(\mathbb{T}^d)$ with the usual symplectic form $(i\cdot,\cdot)_{L^2}$. Therefore, provided that $\ell^p_s(\mathbb{Z}^d) \subset \ell^2(\mathbb{Z}^d)$, a smooth map $\tau : \mathcal{U} \to \ell^p_s(\mathbb{Z}^d)$ is \emph{symplectic} if
\[
\forall u \in \mathcal{U}, \forall v,w \in \ell^p_s(\mathbb{Z}^d), \quad (iv,w)_{L^2} = (i\mathrm{d}\tau(u)(v),\mathrm{d}\tau(u)(w))_{L^2}.
\]
Moreover, if $H,K : \mathcal{U} \to \mathbb{R}$ are two smooth functions such that $\nabla H$ (or $\nabla K$) is $ \ell^p_s(\mathbb{Z}^d)$ valued then the \emph{Poisson bracket} of $H$ and $K$ is defined by
\[
\{ H,K\}(u):= (i \nabla H(u),\nabla K(u))_{L^2}.
\]
Note that, as~usual, we~have
\begin{equation}
\label{eq:pois-brak}
\{ H,K\}
=2i \sum_{k\in \mathbb{Z}^d} \partial_{\overline{u_k}}H \partial_{u_k} K - \partial_{u_k}H \partial_{\overline{u_k}} K.
\end{equation}
Note that this relation allows to extend the Poisson bracket to the case where~$H$ and~$K$ are not real valued.

\subsubsection{Subspaces and projections}
Given any subset $\mathcal{A} \subset \mathbb{Z}^d$, we~always consider~$\mathbb{C}^{\mathcal{A}}$ as the subspace of $\mathbb{C}^{\mathbb{Z}^d}$ of the sequences supported on $\mathcal{A}$. We~denote by $\Pi_{\mathcal{A}} : \mathbb{C}^{\mathbb{Z}^d} \to \mathbb{C}^{\mathcal{A}}$ the projection defined by restriction. We~extend implicitly any function~$F$ on $\mathbb{C}^{\mathcal{A}}$ to a function on $\mathbb{C}^{\mathbb{Z}^d}$ by $F=F\circ \Pi_{\mathcal{A}}$. As a consequence, any function on $\ell^2(\mathcal{A};\mathbb{C}):=\ell^2(\mathbb{Z}^d) \cap \mathbb{C}^{\mathcal{A}}$ can be see as a smooth function in $\ell^2(\mathbb{Z}^d)$ and so all the previous definitions make sense. In particular, we~note that if $F : \mathbb{C}^{\mathcal{A}} \to \mathbb{R} $ is $C^1$ then $\nabla F$ is $\mathbb{C}^{\mathcal{A}}$ valued. Given a positive number $N$, we~set $\Pi_{\geq N} := \Pi_{ \{ n \in \mathbb{Z}^d \mid |n| \geq N \} }$, $\Pi_{N} := \Pi_{ \{ n \in \mathbb{Z}^d \mid |n| \leq N \} }$ and $\mathbb{Z}^d_N = \{n \in \mathbb{Z}^d \mid |n| \leq N \}$.

\section{Hamiltonian setting and high super-actions}
\subsection{A first Hamiltonian formalism} We first introduce a set of multi-indices to describe the polynomials:
\begin{definition}
\label{def:multi}
For $q\geq1$, we~define $\mathcal{N}_{2q}$ the set of multi-indices of degree $2q$ with zero momentum
\[
\ts
\mathcal{N}_{2q} = \big\{ \vec{n} = (n_{1},\dots,n_{2q})\in (\Z^d)^{2q}\mid \sum_{i=1}^{2q}(-1)^{i}n_{i}=0\big\}.
\]
Given $\vec{n}\in\mathcal{N}_{2q}$, we~denote the decreasing rearrangement of $(|n_{1}|,\dots,|n_{2q}|)$ by
\[
n_{1}^{\ast}\geq \cdots \geq n_{2q}^{\ast}.
\]
\end{definition}
Now, we~introduce our main class of polynomials.
\begin{definition}[Real homogeneous polynomial]
\label{def:hamb} Let $q\geq 2$ and $\mathcal{A} \subset \mathbb{Z}^d$. The set $\mathcal{H}_{2q}(\mathcal{A})$ of real homogeneous polynomials of degree $2q$ supported on $ \ell^1(\mathcal{A})$ corresponds to the set of functions
\[
Q(u) = \sum_{\vec{n}\in\mathcal{N}_{2q}} Q_{\vec{n}}u_{n_{1}}\overline{u_{n_{2}}}\cdots \overline{u_{n_{2q}}}=:\sum_{\vec{n}\in\mathcal{N}_{2q}}Q_{\vec{n}}u_{\vec{n}},
\]
where $(Q_{\vec{n}})\in\C^{\mathcal{N}_{2q}}$ satisfies
\begin{enumerate}
\item
(symmetry condition)
\[
\text{$Q_{\vec{n}}$ is symmetric in $(n_{1},n_{3},\dots n_{2q-1})$ and in $(n_{2},n_{4},\dots n_{2q})$.}
\]

\item
(reality condition)
\[
\text{ For all $\vec{n}=(n_{1},n_{2},\dots,n_{2q})\in\mathcal{N}_{2q}$,}\quad Q_{(n_{1},n_{2},\dots n_{2q})} = \overline{Q_{(n_{2q},n_{1},\dots n_{2q-1})}}.
\]

\item
(boundedness condition)
\[
\|Q\|_{\infty} := \underset{\vec{n}\in\mathcal{N}_{2q}}{\sup}\ |Q_{\vec{n}}| <+\infty.
\]

\item
(support)
\[
Q_{(n_{1},n_{2},\dots n_{2q})} \neq 0 \implies \vec{n} \in \mathcal{A}^{2q}.
\]
\end{enumerate}
\end{definition}
\begin{remark}
Thanks to the boundedness condition, these polynomials are smooth functions on $\ell^1(\mathbb{Z}^d)$ and the support condition only means that $Q = Q\circ \Pi_{\mathcal{A}}$.
\end{remark}

\begin{remark} Note that this definition is well suited to the Hamiltonian $H$ of \eqref{eq:nls} (defined by \eqref{eq:hamNLS}). More precisely, in~Fourier variables, the Taylor expansion of $H$ writes
\[
H(u) \mathop{=}_{u \to 0} \frac12 \sum_{n\in \mathbb{Z}^d} \lambda_n^2 |u_n|^2 + \frac12 \sum_{q\geq 2} \frac{f^{(q-1)}(0)}{q!} \sum_{\vec{n}\in\mathcal{N}_{2q}} u_{\vec{n}}
\]
or more quantitatively, for all $r\geq 2$,
\begin{equation}
\label{eq:expansion_ham}
\begin{split}
H(u) &= \frac12 \sum_{n\in \mathbb{Z}^d} \lambda_n^2 |u_n|^2 + \frac12 \sum_{q=2}^r \frac{f^{(q-1)}(0)}{q!} \sum_{\vec{n}\in\mathcal{N}_{2q}} u_{\vec{n}}+ R^{(2r+2)} (u) \\ &=:H^{(\leq 2r)}(u) + R^{(2r+2)} (u),
\end{split}
\end{equation}
where $R^{(2r+2)} $ is a remainder term of order $2r+2$ in the sens that for all $u\in B_s(1)$, $s>\sfrac{d}{2}$
\begin{equation}
\label{eq:rem_NLS}
\| \nabla R^{(2r+2)} (u) \|_{h^s} \lesssim_{r,s} \| u \|_{h^s}^{2r+1}.
\end{equation}
\end{remark}

\begin{definition}[Resonant function] Given a multi-index $\vec{n}\in\mathcal{N}_{2q}$, the resonant function $\Omega_{\vec{n}}$ is
\[
\Omega_{\vec{n}} = \sum_{n=1}^{2q}(-1)^{n+1}\lambda_{n}^{2}.
\]
\end{definition}

\begin{definition}[Quasi-resonant Hamiltonian]\label{def:non-kappa} Let $\kappa>0$, $q\geq2$. A~real homogeneous polynomial $Q \in \mathcal{H}_{2q}$
is $\kappa$-\emph{resonant} if for all $\vec{n}\in\mathcal{N}_{2q}$,
\[
|\Omega_{\vec{n}}| > \kappa\implies Q_{\vec{n}} = 0.
\]
\end{definition}
\setcounter{theorem}{7}
\subsection{A Birkhoff normal form theorem and some basic estimates}

In the energy estimates we will need the following estimate, which is a consequence of Cauchy--Schwarz and Young's convolution inequalities:
\begin{lemma}
\label{lem:young}
Let $q\geq1$, $\iota \in \{-1,1\}^{2q}$ and $a^{(1)},\dots, a^{(2q)}\in \ell^1(\Z^d ; \R_+)$. We~have
\[
\sum_{\substack{\vec{n}=(n_1,\dots,n_{2q})\in(\Z^d)^{2q}\\ \iota_1n_1+\iota_2n_2+\dots +\iota_{2q}n_{2q}=0}}a_{n_1}^{(1)}\dots a_{n_{2q}}^{(2q)}\leq\underset{{\sigma\in\mathfrak{S}_{2q}}}\min \| a^{(\sigma(1))} \|_{\ell^2}\,\| a^{(\sigma(2))} \|_{\ell^2} \prod_{i=3}^{2q}\| a^{(\sigma(i))} \|_{\ell^1}.
\]
\end{lemma}
As usual, we~deduce some useful estimate for the vector fields.
\begin{corollary}[Vector field estimates]\label{cor:vecfi} Given $q\geq 2$, $P\in \mathcal{H}_{2q}(\mathbb{Z}^d)$, $s\geq 0$, $P$~defines a smooth function on $h^s \cap \ell^1$, $\nabla P$ is smooth from $h^s \cap \ell^1$ to $h^s \cap \ell^1$ and for all $u\in h^s \cap \ell^1$, we~have
\[
\|\nabla P\|_{h^s}(u) \lesssim_{q,s} \| P \|_{\infty} \|u\|_{h^s} \|u\|_{\ell^1}^{2q-2}.
\]
\end{corollary}
Then, as~usual, we~deduce the local existence of the Hamiltonian flows of the polynomials of which we summarize the properties that will be useful for us in this paper.

\begin{corollary}[Hamiltonian flows]\label{cor:flow} Let $s > \sfrac{d}{2}$, $r\geq 2$, $C\geq 1$, $\mathcal{A} \subset \mathbb{Z}^d$. For all $\kappa \in (0,1)$, there exists $\varepsilon_* \gtrsim_{r,s,C} \sqrt{\kappa}$ such that for all real polynomial $\chi$ of degree smaller than or equal $2r$ of the form
\[
\chi = \chi^{(4)} + \cdots + \chi^{(2r)},
\]
where
\[
\forall j\in \{ 2,\dots,r\}, \quad \chi^{(2j)} \in \mathcal{H}_{2j}(\mathcal{A})\text{ satisfies } \| \chi^{(2j)} \|_{\infty} \leq C\kappa^{-j+1},
\]
there exists a smooth map
\[
\Phi_{\chi} : \left\{
\begin{array}{rcl} [-1,1] \times \Pi_{\mathcal{A}}^{-1}\Pi_{\mathcal{A}}B_s(\varepsilon_*) & \dpl\to & h^s(\mathbb{Z}^d) \\
(t,u) & \dpl\mto & \Phi_\chi^t(u)
\end{array}\right.
\]
such that, for all $u\in \Pi_{\mathcal{A}}^{-1}\Pi_{\mathcal{A}}B_s(\varepsilon_*) $,
\[
i\partial_t \Phi_\chi^t(u) = (\nabla \chi)\circ \Phi_\chi^t (u) \quad \forall t \in [-1,1],
\]
and for all $t\in [-1,1]$, $\Phi_\chi^t$ is symplectic,
\begin{itemize}
\item close to the identity
\[
\| \Phi_\chi^t(u) - u \|_{h^s} \leq \Big(\frac{\| \Pi_{\mathcal{A}}u \|_{h^s}}{ \varepsilon_* } \Big)^2 \| \Pi_{\mathcal{A}} u\|_{h^s},
\]
\item its derivative is not too big, \ie
\begin{equation}
\label{eq:d-varphi}
\| \mathrm{d} \Phi_\chi^t(u) \|_{h^s \to h^s} \leq 2
\end{equation}
\item and provided that $\Phi_\chi^t(u) \in \Pi_{\mathcal{A}}^{-1}\Pi_{\mathcal{A}}B_s(\varepsilon_*)$, we~have $\Phi_\chi^{-t}(\Phi_\chi^t(u)) = u $.
\end{itemize}
\end{corollary}

\begin{remark}
These projections may seem unusual but they are just a consequence of the block diagonal structure of $\Phi_\chi^t$ :
\[
\Phi_\chi^t = \Pi_{\mathcal{A}} (\Phi_\chi^t)_{| h^s(\mathcal{A})} \oplus \mathrm{id}_{h^s(\mathcal{A}^c)}.
\]
\end{remark}

We also recall the following continuity estimates of the Poisson brackets when applied to polynomials.
\begin{lemma} Let $\mathcal{A} \subset \mathbb{Z}^d$ be a subset of $\mathbb{Z}^d$, $P \in \mathcal{H}_{2p}(\mathcal{A})$ a homogeneous real polynomial of degree $2p\geq 4$ supported on $\ell^1(\mathcal{A})$ and $Q \in \mathcal{H}_{2q}(\mathcal{A})$ a real homogeneous polynomial of degree $2q\geq 4$ supported on $\ell^1(\mathcal{A})$ then their Poisson bracket $\{P,Q\} \in \mathcal{H}_{2(p+q-1)}(\mathcal{A})$ is a real homogeneous polynomial of degree $2(p+q-1)$ supported on $\ell^1(\mathcal{A})$ satisfying the bound
\[
\| \{ P,Q\} \|_{\infty} \lesssim pq \|P\|_{\infty} \| Q\|_{\infty}.
\]
\end{lemma}

As usual, we~deduce of these estimates the following Birkhoff normal form theorem.
\begin{theorem}[Quasi-resonant normal form]
\label{thm:res-nf} Let $r\geq2$, $s> \max(\sfrac{d}{2},1)$, $\mathcal{A} \subset \mathbb{Z}^d$, $\kappa \in (0,1)$ and $P = P^{(4)} + \cdots + P^{(2r)}$ be a real Hamiltonian of degree $2r$ supported on $\ell^1(\mathcal{A})$ with $P^{(2j)} \in \mathcal{H}_{2j}(\mathcal{A})$ for all $j\in \{ 2,\dots,r \}$ and set
\[
H = Z_2^{\mathcal{A}} + P \quad \mathrm{where } \quad Z_2^{\mathcal{A}}(u) = \frac12\sum_{n\in \mathcal{A}} \lambda_n^2 |u_n|^2.
\]
There exist some polynomials $\chi^{(2j)} \in \mathcal{H}_{2j}(\mathcal{A})$, $2\leq j \leq r$, satisfying
\[
\| \chi^{(2j)} \|_{\infty} \lesssim_{ C,r} \kappa^{-j+1},
\]
where $C = \max_{j\geq 2} \kappa^{2-j}\| P^{(2j)} \|_{\infty}$ such that, on $\Pi_{\mathcal{A}}^{-1}\Pi_{\mathcal{A}}B_s(\varepsilon_*) $ (where $\varepsilon_*$ is given by Corollary \ref{cor:flow}),
\[
H\circ \Phi_{\chi}^1 = Z_2^{\mathcal{A}} + Q^{(4)} + \cdots+ Q^{(2r)} + \Upsilon,
\]
where $\chi = \chi^{(4)} + \cdots + \chi^{(2r)}$, the polynomials $Q^{(2j)} \in \mathcal{H}_{2j}(\mathcal{A})$ are some $\kappa$ resonant real homogeneous polynomials satisfying for all $j\in \{2,\dots,r\}$ and all $\vec{n} \in (\mathbb{Z}^d)^4$
\[
\| Q^{(2j)} \|_{\infty} \lesssim_{C,r} \kappa^{-j+2} \quad \mathrm{and} \quad Q^{(4)}_{\vec{n}} = \mathbf{1}_{|\Omega_{\vec{n}}| \leq \kappa} P^{(4)}_{\vec{n}},
\]
and $\Upsilon \in C^1(\Pi_{\mathcal{A}}^{-1}\Pi_{\mathcal{A}}B_s(\varepsilon_*);\mathbb{R})$ is a remainder of order $2(r+1)$ in the sense that for all $u\in \Pi_{\mathcal{A}}^{-1}\Pi_{\mathcal{A}}B_s(\varepsilon_*)$,
\begin{equation}
\label{eq:R-res}
\|\nabla \Upsilon(u)\|_{h^{s}}\lesssim_{C,r,s} \Big(\frac{\| \Pi_{\mathcal{A}} u\|_{h^{s}}}{\sqrt{\kappa}} \Big)^{2r}\|\Pi_{\mathcal{A}} u\|_{h^{s}}.
\end{equation}
\end{theorem}
This is a formulation of Birkhoff normal form theorem, which is by now standard. For the proof, see \cite[Th.\,2.15]{BG22} or \cite[Th.\,2.12]{BGR23}.

\subsection{Almost preservation of the high-super actions}
In this subsection, we~prove Proposition~\ref{prop:high_modes}, in~which we show that up to a bootstrap assumption we control the~$H^s$ norm of the high modes of the solution of \eqref{eq:nls} for very long times. Note that we do not need to use any internal or external parameters.

\begin{proposition}
\label{prop:high_modes} There exist $\delta >0$ and $N_0\geq 2$ depending only on $\mathscr{L}$ such that given $T>0$, $c>1$, $s>d$, $\rho>0$ and $u$ a solution to \eqref{eq:nls} in $C([0,T],h^{s})$ with $\|u(0)\|_{h^{s}} \leq \rho$ and
\begin{equation}
\label{eq:as-ap}
\underset{t\in[0,T]}{\sup}\ \|u(t)\|_{h^{s}}\leq c\rho,
\end{equation}
then for all $N\geq N_0$, for all $r\geq2$, provided that $\varepsilon \lesssim_{c,r,s} 1$ is small enough and $t\leq T$,
\[
\| \Pi_{\geq N} u(t) \|_{h^s}^2 \leq 2^{2s} \| u(0) \|_{h^s}^2 + \rho^3 + TN^{-\Psfrac{\delta}{2}(s-d)} \rho^{3}+T \rho^{2r+1}.
\]
\end{proposition}

The proof is based on Birkhoff normal forms, and on the extension of Bourgain's cluster decomposition lemma \cite{Bou98} to any flat tori, which was proved by Berti and Maspero in \cite[Th.\,2.1]{Berti-Maspero-19}. The use of this decomposition in the context of Birkhoff normal forms was recently initiated by Bambusi, Feola and Montalto in \cite{BFM22}. It~is a way to overcome the small divisor degeneracy discussed in Section \href{https://doi.org/10.5802/jep.300}{1.3.3 of the source article}. In the following, we~will mainly focus on the proof of Proposition~\ref{prop:ap-hi}, from which Proposition~\ref{prop:high_modes} easily follows.

\begin{lemma}[Clustering of eigenvalues of $-\Delta_\mathscr{L}$, {\cite[Th.\,2.1]{Berti-Maspero-19}}]
\label{lem:cluster}
There exist constants $\delta \equiv \delta(d) \in (0,1)$ and $C(\mathscr{L},d)\geq 2$ and a partition of $\Z^d$
\[
\Z^d = \bigcup_{\upsilon }\mathscr{C}_\upsilon,
\]
satisfying the following properties:
\begin{enumerate}
\item The sets $\mathscr{C}_\upsilon$ are finite and, up to a bounded set $\mathscr{C}_{\upsilon_0}$ such that
\[
\underset{n\in \mathscr{C}_{\upsilon_0}}{\max} |n| \leq C(\mathscr{L},d) \quad \mathrm{and} \quad |n| < 2 \implies n\in \mathscr{C}_{\upsilon_0}.
\]
\item They are dyadic in the sense that
\begin{equation}
\label{eq:dydy}
\underset{n\in\mathscr{C}_{{\upsilon}}}{\max} |n| \leq 2\underset{n\in\mathscr{C}_{{\upsilon}}}{\min} |n| =: 2m_\upsilon.
\end{equation}
\item If one denotes $[n_1]$ the class of equivalence of $n_1\in\Z^d$, that is,
\[
[n_1] = \mathscr{C}_\upsilon,\quad \text{where}\ \upsilon\ \text{is such that}\quad n_1\in\mathscr{C}_\upsilon,
\]
then we have the separation property: for all $n_1,n_2\in\Z^d$,
\[
[n_1]\neq[n_2]\implies |n_1-n_2| + |\lambda_{n_{1}}^{2}-\lambda_{n_{2}}^{2}| > (|n_1|+|n_2|)^{\delta(d)}.
\]
\end{enumerate}
\end{lemma}

\begin{remark}
As usual, we~identify the set of the indices $\upsilon$ with the set of the equivalence classes.
\end{remark}

We stress out that this lemma holds for any flat torus, without any Diophantine assumption on the metric. As in \cite{BFM22}, this decomposition allows us to design some almost conserved quantities of \eqref{eq:nls} called \emph{super-action}.

\begin{definition}[Super actions] \label{def:super_actions} Given a mode $n_0\in \mathbb{Z}^d$ the corresponding super-action is defined by
\[
S_{n_0}(u) := \sum_{n\in[n_{0}]} |u_n|^2.
\]
\end{definition}

The cluster being dyadic (see \eqref{eq:dydy}), Proposition~\ref{prop:high_modes} can easily be deduced from the following:
\begin{proposition}[A priori estimates for the high super-actions]\label{prop:ap-hi} Let $T,u,c,s,\rho$ be as in Proposition~\ref{prop:high_modes}. For all $M \geq 2$,
for all $r\geq2$, provided that $\rho\lesssim_r 1$ is small enough and $t\leq T$, we~have
\begin{equation}
\label{eq:ap-hi}
\sum_{m_\upsilon \geq M} m_\upsilon^{2s} \big|S_{\upsilon}(u(t))-S_{\upsilon}(u(0))\big|
\lesssim_{r,s,c} \bigl(\rho^4 + TM^{-\Psfrac{\delta}{2}(s-d)} \rho^{4}+T \rho^{2(r+1)}\bigr),
\end{equation}
where $m_{\upsilon}$ was defined in \eqref{eq:dydy}.
\end{proposition}

\begin{proof}[Proof of Proposition~\ref{prop:high_modes} from Proposition~\ref{prop:ap-hi}] Indeed, to deduce Proposition~\ref{prop:high_modes} from Proposition~\ref{prop:ap-hi}, thanks to the dyadicity of the clusters (see \eqref{eq:dydy}), it is enough to set $N_0 = 3C(\mathscr{L},d)$ and $M=N/2$ to get that
\[
\| \Pi_{\geq N} u \|_{h^s}^2 \leq 2^{2s} \sum_{|n|\geq N} m_{[n]}^{2s} |u_n|^2 \leq 2^{2s} \sum_{ m_\upsilon \geq M} m_\upsilon^{2n} S_\upsilon(u),
\]
and so, thanks to Proposition~\ref{prop:ap-hi}, that
\[
\| \Pi_{\geq N} u(t) \|_{h^s}^2 \leq 2^{2s} \sum_{ m_\upsilon \geq M} m_\upsilon^{2n} S_\upsilon(u(0)) + C_{r,s,c} \, \rho \, \bigl(\rho^3 + TN^{-\Psfrac{\delta}{2}(s-d)} \rho^{3}+T \rho^{2r+1}\bigr),
\]
where $C_{r,s,c}$ is a constant depending only on $(r,s,c)$. We~then conclude the proof of Proposition~\ref{prop:high_modes} by~using the estimate $\sum_{ m_\upsilon \geq M} m_\upsilon^{2s} S_\upsilon(u(0)) \leq \|u(0) \|_{h^s}^2$ and by absorbing the constant $C_{r,s,c} $ thanks to the smallness of $\rho$.
\end{proof}

Before proving Proposition~\ref{prop:ap-hi}, let us consider the following technical lemmas.

\begin{lemma}\label{lem:mu3} Let $\vec{n}\in\mathcal{N}_{2q}$ for some $q\geq2$ be such that $|\Omega_{\vec{n}}|\leq 1$. If there exists a frequency cluster $\mathscr{C}_{\upsilon}$ such that
\[
\{ S_{\upsilon}, u_{\vec{n}}\} \neq 0,
\]
then
\[
n_{1}^{\ast}\geq m_\upsilon \quad \mathrm{and} \quad n_{3}^{\ast}\gtrsim_q (m_\upsilon)^{\delta/2}.
\]
\end{lemma}

\begin{proof}
First, we~note that
\[
\{ S_{\upsilon}, u_{\vec{n}}\} = 2i \biggl(\sum_{j=1}^q \mathbf{1}_{n_{2j} \in \mathscr{C}_\upsilon } - \mathbf{1}_{n_{2j-1} \in \mathscr{C}_\upsilon } \biggr) u_{\vec{n}}.
\]
As a consequence, if $\{ S_{\upsilon}, u_{\vec{n}}\} \neq 0$, there exists $j_0 \in \{1,\dots,2q\}$ and an odd number $\theta\in \mathbb{Z}$ such that $n_{j_0} \in \mathscr{C}_\upsilon$ and $n_{j_0 + \theta} \notin \mathscr{C}_\upsilon$. By~symmetry, without loss of generality, we~assume that $j_0=1$ and $j_0 + \theta= 2$. Since $n_1 \in \mathscr{C}_\upsilon$, we~have $n_{1}^{\ast}\geq m_\upsilon$. Moreover, thanks to the separation property of the clusters, we~have
\[
|n_1-n_2| + |\lambda_{n_{1}}^{2}-\lambda_{n_{2}}^{2}| > (|n_1|+|n_2|)^{\delta} \geq m_\upsilon^{\delta}.
\]
As a consequence, either (by the zero momentum condition)
\[
(2q-2) \max_{j\geq 3} |n_j| \geq |n_1-n_2| \geq \frac{m_\upsilon^{\delta}}2,
\]
or (since $|\Omega_{\vec{n}}|\leq 1$)
\[
(2q-2) \max_{j\geq 3} |n_j|^2 + 1 \gtrsim_{\mathscr{L}} |\lambda_{n_{1}}^{2}-\lambda_{n_{2}}^{2}| \geq \frac{m_\upsilon^{\delta}}2.
\]
In any case, we~have
\[
n_3^* \geq \max_{j\geq 3} | n_j | \gtrsim_{\mathscr{L},q} m_\upsilon^{\delta/2}.\qedhere
\]
\end{proof}

\begin{lemma} Let $s\geq \sigma \geq 0$, $q\geq 2$ and $Q\in\mathcal{H}_{2q}$. If $Q$ is $1$-resonant then for all $u\in h^{s} \cap \ell^1_\sigma$, we~have
\[
\sum_{\upsilon} m_\upsilon^{2s+\sfrac{\sigma \delta}2} |\{S_{\upsilon},Q\}(u)| \lesssim_{q,s} \|Q\|_{\infty} \|u\|_{h^{s}}^{2} \| u \|_{\ell^1_\sigma} \|u\|_{\ell^{1}}^{2q-3}.
\]
\end{lemma}
\begin{proof} First, we~note that if $\upsilon = \upsilon_0$ then $m_\upsilon = 0$ so we do not have to consider this class in the sum. Therefore, we~consider a class $\upsilon \neq \upsilon_0$. Note that it implies that $m_{\upsilon} \geq 2.$ Then, we~note that
\[
\{S_{\upsilon},Q\} = \sum_{\vec{n}\in\mathcal{N}_{2q}} Q_{\vec{n}} \{S_\upsilon, u_{\vec{n}} \}.
\]
Using Lemma~\ref{lem:mu3}, we~deduce that we can reduce the sum to the set of the indices satisfying $n_{1}^{\ast}\geq m_\upsilon \geq 2$ and
\[
n_{3}^{\ast}\gtrsim_{q}m_\upsilon^{\delta/2}.
\]
Moreover, using the zero momentum condition, we~note that $n_{2}^{\ast}\geq (2q)^{-1} n_1^\ast \geq (2q)^{-1} m_\upsilon$. It~follows that
\begin{align*}
\sum_{\upsilon } m_\upsilon^{2s+\sfrac{\sigma \delta}2}|\{S_{\upsilon},Q\}(u)|
&\lesssim_q \sum_{\upsilon} m_\upsilon^{2s+\sfrac{\sigma \delta}2} \sum_{\substack{n_{2}^{\ast} \gtrsim_q m_\upsilon \\ n_{3}^{\ast} \gtrsim_{q} m_\upsilon^{\delta/2} }}|Q_{\vec{n}}||u_{n_1}| \cdots |u_{n_{2q}}| \\
&\lesssim_{q,s} \| Q\|_{\infty} \sum_{\vec{n}\in\mathcal{N}_{2q}} (n_{1}^{\ast})^{s} (n_{2}^{\ast})^{s} (n_{3}^{\ast})^{\delta \sigma /2} |u_{n_{1}}|\cdots |u_{n_{2q}}| \\
&\lesssim_{q,s} \| Q\|_{\infty} \sum_{\varphi \in \mathfrak{S}_{2q}} \sum_{\vec{n}\in\mathcal{N}_{2q}} |n_{\varphi_1}|^{s} |n_{\varphi_2}|^{s} |n_{\varphi_3}|^{\delta \sigma /2} |u_{n_{1}}|\cdots |u_{n_{2q}}|.
\end{align*}
And so we get the expected estimate by applying the Young estimate of Lemma~\ref{lem:young}.
\end{proof}

By applying this lemma with $\sigma = s-d$ and by considering only the classes such that $m_\upsilon$ is large enough we deduce the following estimate.
\begin{corollary}
\label{cor:super}Let $s\geq d$, $q\geq 2$ and $Q\in\mathcal{H}_{2q}$. If $Q$ is $1$-resonant, then, for all $u\in h^{s}$ and $M\geq 2$, we have
\begin{equation}
\label{eq:super}
\sum_{m_\upsilon \geq M} m_\upsilon^{2s} |\{S_{\upsilon},Q\}(u)| \lesssim_{q,s} \|Q\|_{\infty} M^{-\Psfrac{\delta}{2}(s-d)}\|u\|_{h^{s}}^{3} \|u\|_{\ell^{1}}^{2q-3}.
\end{equation}
\end{corollary}

Now we start proving Proposition~\ref{prop:ap-hi} about the preservation of the super-actions of \eqref{eq:nls}.

\begin{proof}[Proof of Proposition~\ref{prop:ap-hi}]
Recall the Hamiltonian formulation \eqref{eq:hamNLS} of \eqref{eq:nls} and the Taylor expansion $H = H^{(\leq 2r)} + R^{(2r+2)}$ of its Hamiltonian \eqref{eq:expansion_ham}. We apply the Birkhoff normal form theorem~\ref{thm:res-nf} with $\kappa=1$ and $\mathcal{A}= \mathbb{Z}^d$ to the Hamiltonian $H^{(\leq 2r)}$ of degree $2r$ to put it in $1$-resonant normal form up to order $2r$. This gives an auxiliary Hamiltonian $\chi = \chi^{(4)} + \cdots + \chi^{(2r)}$ with $\| \chi^{(2j)}\|_\infty \lesssim_r 1$ and, through Corollary \ref{cor:flow}, a radius $\varepsilon_\ast \gtrsim_{r,s} 1$ such that on $B_s(\varepsilon_*)$
\[
\widetilde{H}:=H^{(\leq 2r)}\circ\Phi_\chi^1 =Z_2+ \sum_{q=2}^{r}\widetilde{Q}^{(2q)}+\Upsilon =: \widetilde{H}^{(\leq 2r)} + \Upsilon,
\]
where $\widetilde{Q}^{(2q)} \in \mathcal{H}_{2q}$ are 1-resonant homogeneous polynomials satisfying $\| Q^{(2q)}\|_{\infty} \lesssim_q 1 $, and
the remainder $\Upsilon$ is of order $2(r+1)$ in the sense that for all $w\in B_s(\varepsilon_*)$
\begin{equation}
\label{eq:bientotnoel}
\|\nabla \Upsilon(w)\|_{h^{s}}\lesssim_{r,s} \| w\|_{h^s}^{2r+1}.
\end{equation}
From now on, we assume that the solution $u$ of \eqref{eq:nls} satisfies \hbox{$\| u(0) \|_{h^s} \!\!\leq\! \rho \!\leq\! 10^{-2}\! c^{-1}\!\varepsilon_\ast$}. It~follows that for all $t\in [0,T]$, we~have $\|u(t) \|_{h^s} \leq c \rho\leq 10^{-2} \varepsilon_\ast$ and so that $v = \Phi_\chi^{-1}(u) \in C^0([0,T];h^s)$ is well-defined for $t\in [0,T]$. Moreover, $ \Phi_\chi^{-1}$ being close to the identity, we~have
\begin{equation}
\label{eq:pr-var}
\|u(t)-v(t)\|_{h^{s}}\leq \varepsilon_{\ast}^{-2}\|u(t)\|_{h^{s}}^{3} \lesssim_{r,c} \rho^3,\quad \text{and \ so}\quad \|v(t)\|_{h^{s}}\lesssim_{r,c} \rho.
\end{equation}
Without loss of generality, we~can assume that\footnote{It can be done by approximating any $u(0)$ by smoother initial data (\eg trigonometric polynomials) and using the preservation of the regularity for \eqref{eq:nls} (\ie that while the $\ell^1$ norm of the solution is bounded, its $h^{s+2}$ norm does not blow up). Such an approximation is classical and is done in details, for example, in~\cite{BG22}.} $u(0) \in h^{s+2}$, so that $u\in C^1([0,T];h^s)$ and $v \in C^1([0,T];h^s)$. Since $\Phi_\chi^{-1}$ is symplectic, using the chain rule, we~get that
\[
i\partial_{t}v = \nabla \widetilde{H}^{(\leq 2q)}(v) + \nabla \Upsilon (v) + \mathrm{d}\Phi_\chi^{-1}(u)\bigl(\nabla R^{(2r+2)}(u)\bigr).
\]
As a consequence, one can bound the variation of the super-actions as follows:
\begin{multline}
\label{eq:pr-v}
\sum_{m_\upsilon \geq M}m_\upsilon^{2s}\Big|\frac{d}{dt}S_{\upsilon}(v(t))\Big|
\leq \sum_{m_\upsilon \geq M}m_\upsilon^{2s}|\{S_{\upsilon},\widetilde{H}^{(\leq 2r)}\}(v)|
+ \|\nabla \Upsilon (v)\|_{h^{s}}\|v\|_{h^{s}}\\+\|\mathrm{d}\Phi_\chi^{-1}(u)\|_{h^{s}\to h^{s}}\|\nabla R^{(2r+2)}(u)\|_{h^{s}}\|v\|_{h^{s}}.
\end{multline}
According to Corollary \ref{cor:super} the first term on the right-hand side is bounded (up to a constant depending only on $r$ and $s$) by
\[
\underset{2\leq q \leq r}{\max} \|\widetilde{Q}^{(2q)}\|_{\infty} M^{-\Psfrac{\delta}{2}(s-d)}\|v(t)\|_{h^{s}}^{3} \|v(t)\|_{\ell^{1}}^{2q-3} \lesssim_{r,s,c} M^{-\Psfrac{\delta}{2}(s-d)} \rho^4.
\]
Moreover, the estimate \eqref{eq:d-varphi} on $\mathrm{d}\Phi_\chi^{t}$ together with the estimate \eqref{eq:rem_NLS} on $\nabla R^{(2r+2)}$ yield
\[
\|\mathrm{d}\Phi_\chi^{-1}(u)\|_{h^{s}\to h^{s}}\| \nabla R^{(2r+2)} (u)\|_{h^{s}} \lesssim_{r,s,c} \rho^{2r+1}.
\]
Using also the estimate \eqref{eq:bientotnoel} on $\nabla \Upsilon$,
we conclude that
\[
\sum_{m_\upsilon \geq M}m_\upsilon^{2s}\Big|\frac{d}{dt}S_{\upsilon}(v(t))\Big| \lesssim_{c,r,s} M^{-\Psfrac{\delta}{2}(s-d)} \rho^4+ \rho^{2r+1}.
\]
Finally, using the fundamental theorem of calculus and that $v$ is close to $u$ (see~\eqref{eq:pr-var}, it provides the term $\rho^4$ without factor $T$ in \eqref{eq:ap-hi}), and we~get the expected estimate.
\end{proof}
\begin{thebibliography}{BGR23}
\bibitem[BG25]{BG22} J. Bernier \& B. Grébert, “Almost global existence for some nonlinear Schrödinger equations on Td in low regularity”, Ann. Inst. Fourier (Grenoble) (2025), Online first.
\bibitem[BGR23]{BGR23} J. Bernier, B. Grébert \& T. Robert – “Dynamics of quintic nonlinear Schrödinger equations in H 2/5+ (T)”, 2023, To appear in Transactions of the AMS, arXiv:2305.05236.
\bibitem[Bou98]{Bou98} J. Bourgain – “Quasi-periodic solutions of Hamiltonian perturbations of 2D linear Schrödinger equations”, Ann. of Math. (2) 148 (1998), no. 2, p. 363–439.
\bibitem[BM19]{Berti-Maspero-19} M. Berti \& A. Maspero – “Long time dynamics of Schrödinger and wave equations on flat tori”, J. Differential Equations 267 (2019), no. 2, p. 1167–1200.
\bibitem[BFM24]{BFM22} D. Bambusi, R. Feola \& R. Montalto – “Almost global existence for some Hamiltonian PDEs with small Cauchy data on general tori”, Comm. Math. Phys. 405 (2024), no. 1, article no. 15 (50 pages).
\end{thebibliography}
\end{document}
